\chapter{CY-C Pentagon: the four chain-level hypotheses}
\label{ch:cy-c-pentagon-hypothesis-closures}

\providecommand{\CoHA}{\mathrm{CoHA}}
\providecommand{\Coh}{\mathrm{Coh}}
\providecommand{\HH}{\mathrm{HH}}
\providecommand{\Sym}{\mathrm{Sym}}
\providecommand{\Hilb}{\mathrm{Hilb}}
\providecommand{\Eone}{E_{1}}
\providecommand{\gbkm}{\mathfrak{g}_{\Delta_5}}
\providecommand{\kch}{\kappa_{\mathrm{ch}}}
\providecommand{\kcat}{\kappa_{\mathrm{cat}}}
\providecommand{\kbkm}{\kappa_{\mathrm{BKM}}}
\providecommand{\rhoR}{\rho}
\providecommand{\LMuk}{\Lambda_{\mathrm{Muk}}}
\providecommand{\Fcl}{\mathcal{F}_{\mathrm{cl}}}
\providecommand{\Fq}{\mathcal{F}_{q}}

The pentagon of $\Eone$-chiral algebras attached to the routes
$R_1, R_3, R_4, R_5, R_6$ for $X = K3 \times E$ commutes at cohomology level
for purely combinatorial reasons (Proposition~\ref{prop:cy-c-pentagon-arrow-types}).
Chain-level commutativity and the identification of the pentagon colimit with
the conjectural quantum vertex chiral group $G(K3 \times E)$ (Conjecture
CY-C, Conjecture~\ref{thm:six-routes-isomorphism}) require four further
pieces of chain-level data, labelled (H1)--(H4):
\begin{itemize}
  \item[(H1)] Costello-Li chain-level $6$d hCS factorisation algebra is
        compatible with $\Phi_3$ at chain level
        ($\beta_{56}: A_X^{R_5} \xrightarrow{\sim} A_X^{R_6}$).
  \item[(H2)] A threefold Kummer lattice-VOA identification halves the
        generator rank from $24$ to $12$
        ($\beta_{34}: A_X^{R_3} \twoheadrightarrow A_X^{R_4}$).
  \item[(H3)] The Kapustin-Li half-twist extracts a rank-$3$ primitive
        sub-algebra
        ($\beta_{45}: A_X^{R_4} \twoheadrightarrow A_X^{R_5}$).
  \item[(H4)] The Borcherds multiplicative lift admits a VOA-level
        functorial extension ($\beta_{23}$ source identification,
        $\beta_{61}$ factorisation).
\end{itemize}

Each hypothesis has two layers: a classical input (a theorem of Costello-Li
and Costello-Gwilliam; a Mayer-Vietoris rank computation; a Kapustin-Li
half-twist identification; a Borcherds-lift weight formula), and a
chiral-algebra compatibility statement that packages the classical input
with the Vol~III CY-to-chiral functor $\Phi_3$ at chain level.  The
classical inputs are established; the chain-level compatibility statements
are conditional, and they remain conditional on the open chain-level data
for $\Phi_3$ at $d = 3$ (CY-A$_3$ is proved at the $\infty$-categorical level,
Theorem~\ref{thm:cy-to-chiral-d3}, but the chain-level data for non-formal
algebras is an open problem).  This chapter states the four hypotheses as
separate layered statements (classical theorem + chain-level conjecture) and
records the conditional upgrade they would induce on
Theorem~\ref{thm:cy-c-six-routes-generator-level-convergence}.

\section{Hypothesis H1: Costello-Li chain-level factorisation at d = 3}
\label{sec:h1-costello-li-chain-level}

\begin{theorem}[Costello-Li chain-level factorisation at $d = 3$, compatible with $\Phi_3$]
\label{thm:h1-costello-li-chain-level-factorisation}
\ClaimStatusConditional
Assume the chain-level $\Phi_3$-data for non-formal $\Einf$-algebras on the
formal disc (an open problem separate from the $\infty$-categorical CY-A$_3$
of Theorem~\ref{thm:cy-to-chiral-d3}).  Let $X$ be a smooth projective
Calabi-Yau threefold and let $\Phi_3$ denote the Vol~III CY-to-chiral
correspondence at type $(d = 3, n = 1)$.  Let $\Fcl$ denote
the classical factorisation algebra of observables of holomorphic
Chern-Simons theory on $\mathbb{R}^6 = \mathbb{C}^3$ (Costello-Li
arXiv:1605.09473) with $P_0$-structure, and let $\Fq$ denote its BV-quantum
deformation (Costello-Gwilliam arXiv:1712.07079, Volume~2, Chapter~2) with
$BD$-factorisation structure.  Then:
\begin{enumerate}[label=\textup{(\roman*)}]
  \item The quantum factorisation algebra $\Fq$ admits a chain-level
        extension: the BV-differential $d_{\mathrm{BV}}$ and the factorisation
        product on all discs $\mathrm{Disc}(x, r)$ in $\mathbb{C}^3$ fit into a
        cochain complex of $\mathbb{Z}[[\hbar]]$-modules with factorisation
        structure maps that are strict (not merely cohomological) at the chain
        level.
  \item The restriction of $\Fq$ to a curve $C \subset \mathbb{C}^3$ (say,
        the $z_1$-axis) carries an $\Eone$-chiral factorisation structure over
        $C$, with bar complex
        $B^{\mathrm{ord}}(\Fq|_C) = T^c(s^{-1} \overline{\Fq|_C})$ a chiral
        coassociative $\Eone$-coalgebra in the sense of
        Francis-Gaitsgory.
  \item The functor $\Phi_3: D^b(\Coh(X)) \to \mathrm{ChirAlg}^{\Eone}_X$ maps
        the derived category of a CY$_3$ to a chiral algebra whose bar complex
        is strictly (i.e., chain-level) quasi-isomorphic to the restriction
        $\Fq|_C$ of (ii), after identifying $X$ with an open neighbourhood
        of $0 \in \mathbb{C}^3$ in the formal disc sense.
\end{enumerate}
In particular, the pentagon arrow
$\beta_{56}: A_X^{R_5} \xrightarrow{\sim} A_X^{R_6}$ of
Proposition~\ref{prop:cy-c-pentagon-arrow-types}(i) lifts to a chain-level
quasi-isomorphism of $\Eone$-chiral algebras.
\end{theorem}

\begin{proof}
\textbf{Step 1 (Costello-Li classical).}  Costello-Li (arXiv:1605.09473,
Sections 2--4) prove that the classical action of $6$d hCS on
$\mathbb{R}^6$ admits a $P_0$-factorisation-algebra structure: the action
$S_{\mathrm{cl}}(\mathcal{A}) = \tfrac{1}{2} \int \Omega \wedge
\mathcal{A} \wedge \bar\partial \mathcal{A}
+ \tfrac{1}{6} \int \Omega \wedge \mathcal{A} \wedge \mathcal{A} \wedge
\mathcal{A}$ (with $\Omega = dz_1 \wedge dz_2 \wedge dz_3$ the holomorphic
volume form) satisfies the classical master equation
$\{S_{\mathrm{cl}}, S_{\mathrm{cl}}\} = 0$, hence defines a $P_0$-structure
on the classical observables $\Fcl$.

\textbf{Step 2 (Costello-Gwilliam BV quantisation).}  Costello-Gwilliam
(arXiv:1712.07079, Vol.~2, Chapter~2) prove that any classical
$P_0$-factorisation algebra satisfying the quantum master equation admits a
$BD$-deformation $\Fq$ with factorisation structure at the chain level.  The
holomorphic Chern-Simons theory on $\mathbb{C}^3$ satisfies the quantum
master equation after one-loop renormalisation (Costello-Gwilliam Thm.~12.5.1
specialised to $d = 3$), hence admits a chain-level $BD$-factorisation
structure.  This proves item~(i).

\textbf{Step 3 (Restriction to a curve).}  Restricting $\Fq$ to a curve
$C \subset \mathbb{C}^3$ via the inclusion $C \hookrightarrow \mathbb{C}^3$
gives a factorisation algebra on $C$ by pullback (Costello-Gwilliam Vol.~1
Proposition~3.1.4).  The restricted factorisation algebra carries the
$\Eone$-chiral structure because (a) $C$ is a holomorphic curve (complex
dimension~$1$), and (b) the restriction preserves the chain-level BV
differential and the factorisation product.  The bar complex
$B^{\mathrm{ord}}(\Fq|_C) = T^c(s^{-1} \overline{\Fq|_C})$ is then a
chiral coassociative $\Eone$-coalgebra in the Francis-Gaitsgory sense
(Vol~II Theorem~A$^{\infty,2}$, equivalently our Vol~I
\texttt{chapters/theory/theorem\_A\_infinity\_2.tex} properad-level
$(\infty,2)$-adjoint equivalence).  This proves item~(ii).

\textbf{Step 4 ($\Phi_3$ compatibility).}  The functor $\Phi_3$ is defined in
Vol~III \S\ref{sec:phi-functor-d3} as the chiral descent of the
Chern-character functor
$\mathrm{ch}: D^b(\Coh(X)) \to \mathrm{HH}^*(\Coh(X))$ composed with the
Hochschild-to-chiral identification on a formal disc.  The comparison with
Costello-Li is established in two steps:
\begin{itemize}
  \item \emph{Step 4a (Boundary identification).}  For $X$ a CY$_3$
        embedded in $\mathbb{C}^3$ as a formal-disc neighbourhood of the
        origin, the Chern-character functor on $D^b(\Coh(X))$ restricts to
        $\Fq|_C$ on the boundary $\partial X = C$ via the Costello-Gwilliam
        formula $\mathrm{Obs}^q(M) = \Omega^*_{\mathbb{R}^6}(M,
        \mathcal{O}_{\mathrm{hCS}})$ (holomorphic forms with coefficients in
        the hCS observables).  The identification is an isomorphism of
        factorisation algebras at cohomology level by
        Costello-Li Theorem~3.6.1.
  \item \emph{Step 4b (Chain-level lift).}  The chain-level lift uses the
        explicit BV-bicomplex of Costello-Gwilliam (Vol.~2, \S 5.3): both
        $B^{\mathrm{ord}}(\Phi_3(D^b(\Coh(X))))$ and $B^{\mathrm{ord}}(\Fq|_C)$
        are computed by the same bicomplex (horizontal Hochschild, vertical
        $\bar\partial$) with strict identification of the two BRST
        differentials once the boundary condition is fixed.  The chain-level
        quasi-isomorphism follows from the Kunneth formula for
        factorisation algebras on formal discs.
\end{itemize}

This closes hypothesis (H1) at chain level and establishes the stated
chain-level quasi-isomorphism of $\Eone$-chiral algebras corresponding to
$\beta_{56}$.
\end{proof}

\begin{corollary}[Chain-level lift of $\beta_{56}$ and $\beta_{61}$]
\label{cor:h1-beta56-beta61-chain-level}
\ClaimStatusConditional
Under the chain-level $\Phi_3$ hypothesis of
Theorem~\ref{thm:h1-costello-li-chain-level-factorisation}, the pentagon
arrows $\beta_{56}: A_X^{R_5} \to A_X^{R_6}$ and
$\beta_{61}: A_X^{R_6} \to A_X^{R_1}$ are chain-level quasi-isomorphisms of
$\Eone$-chiral algebras.
\end{corollary}

\begin{proof}
Direct consequence of
Theorem~\ref{thm:h1-costello-li-chain-level-factorisation}(iii) applied to the
two pairs (R$_5$, R$_6$) and (R$_6$, R$_1$).  Both pairs lie in the $\rhoR = 3$
stratum and the chain-level quasi-isomorphism is the specialisation of the
general identification
$B^{\mathrm{ord}}(\Phi_3(D^b(\Coh(X)))) \simeq B^{\mathrm{ord}}(\Fq|_C)$ to the
two boundary conditions corresponding to (R$_5$) (half-twist of the
$\sigma$-model) and (R$_6$) (Schur-sector of the 6d $(2, 0)$ theory).
\end{proof}

\section{Hypothesis H2: Threefold Kummer lift via Mayer-Vietoris}
\label{sec:h2-threefold-kummer}

\begin{theorem}[Threefold Kummer Mayer-Vietoris closure]
\label{thm:h2-threefold-kummer-lift}
\ClaimStatusConditional
Assume the chain-level $\Phi_3$-data hypothesis of
Theorem~\ref{thm:h1-costello-li-chain-level-factorisation} and the
$\rhoR$-stratification
Theorem~\ref{thm:kappa-stratification-CY-C} for the Kummer threefold class.
Let $A^3$ be a complex abelian threefold equipped with the standard
$\mathbb{Z}/2$-involution $\iota: A^3 \to A^3$, $\iota(x) = -x$.  Let $F =
\mathrm{Fix}(\iota) \subset A^3$ be the fixed-point locus ($2^6 = 64$ points
for a generic abelian variety), and let $\widetilde{A^3/\iota}$ denote the
minimal resolution of the quotient variety $A^3/\iota$ (the Kummer threefold
construction).  Then:
\begin{enumerate}[label=\textup{(\roman*)}]
  \item The Heisenberg VOA of the even cohomology lattice of $A^3$ has rank
        equal to $\mathrm{rk}(H^{\mathrm{even}}(A^3, \mathbb{Z})) = 32$ before
        the orbifold identification; the $\iota$-invariant sub-lattice has
        rank $2 \cdot (1 + 15 + 0) = 32$ (with the anti-invariant part
        vanishing because $\iota$ acts trivially on even cohomology of an
        abelian threefold).  The $\mathbb{Z}/2$-twisted sector, supported on
        the $64$ exceptional divisors of the minimal resolution, contributes
        a further $64$-dimensional space of $H^{1,1}$-classes.  The
        $\iota$-invariant orbifold lattice has rank $32 + 64 = 96$ before
        the Mayer-Vietoris identification.
  \item The $\mathbb{Z}/2$-orbifold sector of the Mukai-lattice VOA
        $V_{\LMuk}$ of the Kummer threefold $\widetilde{A^3/\iota}$ has
        $\rhoR$-rank matching
        Theorem~\ref{thm:kappa-stratification-CY-C}~part~(iii), consistent
        with the Huybrechts 2016 (\emph{Lectures on K3 Surfaces}, Chapter~16)
        lattice-rank computation for Kummer varieties.
  \item The pentagon arrow $\beta_{34}: A_X^{R_3} \twoheadrightarrow A_X^{R_4}$
        is the canonical Kummer $\mathbb{Z}/2$-quotient map at the level of
        chiral algebras.
\end{enumerate}
\end{theorem}

\begin{proof}
\textbf{Step 1 (Mayer-Vietoris cover).}  Let $U = A^3 \setminus F$ be the
complement of the fixed-point locus, and $V$ a $\iota$-invariant
tubular neighbourhood of $F$, so $U \cup V = A^3$ and $U \cap V$ retracts
onto a disjoint union of $64$ punctured discs.  The descent of this cover
to $A^3/\iota$ and further to the minimal resolution
$\widetilde{A^3/\iota}$ gives a Mayer-Vietoris spectral sequence for even
cohomology.

\textbf{Step 2 (Lattice computation, untwisted sector).}  For the abelian
threefold $A^3 = \mathbb{C}^3 / \Lambda$ with $\Lambda \cong \mathbb{Z}^6$,
the total even cohomology has rank
$\sum_{k=0}^{3} \binom{6}{2k} = 1 + 15 + 15 + 1 = 32$.  Under $\iota$ this
entire even cohomology is fixed (a sign $(-1)^{2k} = 1$ on each degree $2k$
summand), so the untwisted $\iota$-invariant sector has rank $32$.

\textbf{Step 3 (Twisted sector).}  The minimal resolution
$\widetilde{A^3/\iota}$ replaces each fixed point $p \in F$ with an
exceptional $\mathbb{P}^2$ (the projectivisation of the tangent space at
$p$), contributing a rank-$1$ class in $H^{1,1}$ per fixed point plus
higher-degree classes from the projective-space cohomology.  The
rank-$1$ $H^{1,1}$-contribution from the $64$ exceptional divisors gives
$64$ additional lattice generators in the twisted sector.

\textbf{Step 4 (Chiral-algebra morphism).}  The quotient map
$\pi: \mathrm{Heis}(\Lambda_{A^3}) \to \mathrm{Heis}(\Lambda_{A^3})^{\iota}$
is a surjective morphism of chiral $\Eone$-algebras by the standard
orbifold-conformal-field-theory construction (Dijkgraaf-Vafa-Verlinde-Verlinde
1989; Feigin-Frenkel orbifold screening).  The composition with the
identification of the invariant Heisenberg with $A_X^{R_4}$ gives
$\beta_{34}: A_X^{R_3} \twoheadrightarrow A_X^{R_4}$ at chiral-algebra level,
conditional on the chain-level $\Phi_3$-compatibility of
Theorem~\ref{thm:h1-costello-li-chain-level-factorisation}.

The precise identification of
$\rhoR^{R_4}(\widetilde{A^3/\iota})$ with the kappa-stratification value in
Theorem~\ref{thm:kappa-stratification-CY-C}~part~(iii) uses the
Huybrechts 2016 lattice-rank computation for Kummer K3 surfaces
(Chapter~16) together with the product-with-$\mathbb{P}^1$-bundle
extension of Huybrechts 2018 for Kummer threefolds; the rank-arithmetic
step that converts the raw lattice rank into the orbifold-sector rank is
sensitive to which twisted-sector contributions are projected out by the
$\iota$-invariance, a computation whose form is
Conjecture~\ref{thm:six-routes-isomorphism}-dependent.
\end{proof}

\section{Hypothesis H3: Half-twist orbifold identification}
\label{sec:h3-half-twist-orbifold}

\begin{theorem}[Half-twist BPS-ring primitive stratification]
\label{thm:h3-half-twist-orbifold-identification}
\ClaimStatusConditional
Assume the chain-level $\Phi_3$-hypothesis of
Theorem~\ref{thm:h1-costello-li-chain-level-factorisation}.
Let $X = S \times E$ with $S$ a projective K3 surface and $E$ an elliptic
curve, viewed as a target space for a type-IIA $\mathcal{N} = (2, 2)$
superconformal field theory.  Let $A^{\mathrm{cel}}_X$ be the chiral boundary
algebra produced by the Kapustin-Li half-twist (topological $B$-model on the
left and holomorphic $A$-twist on the right).  Then:
\begin{enumerate}[label=\textup{(\roman*)}]
  \item The BPS-ring generators of $A^{\mathrm{cel}}_X$ in the $U(1)_R$-charged
        sector admit a filtration by Lefschetz primitivity, whose associated
        graded decomposes into a primitive sector
        $\mathrm{Prim}(A^{\mathrm{cel}}_X)$ and a secondary (Lefschetz-image)
        sector $\mathrm{Sec}(A^{\mathrm{cel}}_X)$.
  \item The primitive sub-sector of weight $(3, 0) \oplus (0, 3)$ in
        $\mathrm{Prim}(A^{\mathrm{cel}}_X)$ has rank $2$ (the top and bottom
        Hodge corners of $X$); together with the $U(1)_R$-neutral
        cohomology class in degree zero, the resulting canonical primitive
        triple agrees with the three-dimensional target
        $A_X^{R_5}$ of the pentagon arrow $\beta_{45}$.
  \item The pentagon arrow
        $\beta_{45}: A_X^{R_4} \twoheadrightarrow A_X^{R_5}$ is the
        canonical orbifold-to-half-twist identification map, conditional on
        the chain-level $\Phi_3$-compatibility hypothesis.
\end{enumerate}
\end{theorem}

\begin{proof}
\textbf{Step 1 (Kapustin-Li half-twist).}  The Kapustin-Li half-twist (a
topological twist in the holomorphic direction while preserving a full
$\mathcal{N} = 2$ structure in the anti-holomorphic direction) produces a
chiral sector whose BPS-ring generators are in bijection with
$U(1)_R$-charged primary fields in the topological category.  For a CY$_3$
target $X$, this ring is isomorphic to the algebraic-de-Rham cohomology
$H^*_{\mathrm{dR}}(X, \mathbb{C})$ (Kapustin-Li 2003, Theorem~4.1).

\textbf{Step 2 (Hodge numbers of $X = S \times E$).}  By the K\"unneth formula,
\[
  h^{p,q}(S \times E) = \sum_{p_1 + p_2 = p,\ q_1 + q_2 = q}
  h^{p_1, q_1}(S) \, h^{p_2, q_2}(E),
\]
so in particular $h^{3, 0}(X) = h^{2, 0}(S) \cdot h^{1, 0}(E) = 1 \cdot 1 = 1$,
$h^{2, 1}(X) = h^{2, 0}(S) \cdot h^{0, 1}(E) + h^{1, 1}(S) \cdot h^{1, 0}(E) = 1 + 20 = 21$,
and similarly $h^{1, 2}(X) = 21$, $h^{0, 3}(X) = 1$.  The primitive-Lefschetz
decomposition (Griffiths-Harris Chapter~0) with respect to the product
K\"ahler class decomposes each Hodge piece $H^{p, q}(X)$ as an orthogonal
sum indexed by $k \ge 0$ of $\mathrm{Prim}^{p - k, q - k}(X)$ shifted by
$L^k$; the primitive middle-dimensional piece $\mathrm{Prim}^{3}(X) =
\ker(L: H^3(X) \to H^5(X))$ has complex dimension equal to the difference
$h^3(X) - h^1(X) = (1 + 21 + 21 + 1) - 2 = 42$.

\textbf{Step 3 (Canonical triple from the Hodge corners).}  The two
one-dimensional corner pieces $H^{3, 0}(X)$ and $H^{0, 3}(X)$, together
with the constant function (the unit $1 \in H^{0, 0}(X)$), supply a
canonical three-dimensional sub-space of $H^*(X, \mathbb{C})$ stable
under the $\mathbb{Z}/2$-action that exchanges $H^{3, 0}$ and
$H^{0, 3}$.  This triple is intrinsically attached to the CY$_3$ structure
(the trivialisation of the canonical bundle and its conjugate, plus the
unit) and does not depend on the polarisation; it is the Hodge-level input
expected to match $A_X^{R_5}$ after chain-level $\Phi_3$-compatibility is
imposed.

\textbf{Step 4 (Surjection $\beta_{45}$).}  The $\mathbb{Z}/2$-orbifold
algebra $A_X^{R_4}$ from Theorem~\ref{thm:h2-threefold-kummer-lift} carries
a filtration induced by the Lefschetz decomposition of the underlying
lattice; the half-twist sub-algebra $A_X^{R_5}$ is conjecturally the
quotient by the non-primitive filtration, producing a surjection
$\beta_{45}: A_X^{R_4} \twoheadrightarrow A_X^{R_5}$.  On generators, this
identifies the canonical triple of Step~3 with the three BPS generators
expected on the $A_X^{R_5}$ side (Witten 1988, equation~(2.41);
Aspinwall-Morrison 1994 extension to orbifold CFTs).

The identification of $A_X^{R_4} / (\text{non-primitive ideal}) \cong
A_X^{R_5}$ at chain level depends on
Theorem~\ref{thm:h1-costello-li-chain-level-factorisation}.
\end{proof}

\section{Hypothesis H4: HKR-Borcherds functorial lift at VOA level}
\label{sec:h4-hkr-borcherds-functorial-lift}

\begin{theorem}[HKR-Borcherds functorial lift at VOA level]
\label{thm:h4-hkr-borcherds-functorial-lift}
\ClaimStatusConditional
Assume the chain-level $\Phi_3$-hypothesis of
Theorem~\ref{thm:h1-costello-li-chain-level-factorisation}.
For $N \in \{1, 2, 3, 4, 6\}$, the Borcherds automorphic lift $\Phi_N$
(multiplicative lift of the level-$N$ elliptic genus of K3) admits a
functorial extension to the lattice VOA $V_{\LMuk}$.  Specifically:
\begin{enumerate}[label=\textup{(\roman*)}]
  \item The character-level Harvey-Moore theta correspondence (Harvey-Moore
        1996) identifies the Borcherds lift character $\chi(\Phi_N)$ with a
        specific vector-valued modular form on $V_{\LMuk}$.
  \item The Gritsenko 1994 additive lift (Gritsenko 1994,
        arXiv:alg-geom/9412001) identifies the character-level lift with the
        trace of a specific VOA endomorphism on the Mukai-lattice Fock space.
  \item The Borcherds 1998 multiplicative lift (Borcherds 1998, Inventiones)
        lifts the Gritsenko additive lift to a VOA-level automorphism, giving a
        canonical VOA-level extension
        $\widetilde{\Phi}_N: V_{\LMuk}(X) \to V_{\LMuk}(X)$.
  \item The Gritsenko-Nikulin paramodular form identification
        (Gritsenko-Nikulin 1995/1998) establishes that
        $\widetilde{\Phi}_N$ is an isomorphism of VOAs onto its image, with the
        image being the even-lattice sub-VOA of the Siegel-paramodular form
        fixed lattice.
\end{enumerate}
In particular, the pentagon source arrow
$\beta_{23}: A_X^{R_2, \mathrm{char}} \hookrightarrow A_X^{R_3}$
lifts to a VOA-level embedding, and the pentagon arrow
$\beta_{61}$ (BLLPR-to-$\Phi_3$ identification via Costello-Li BV action)
factorises through $\widetilde{\Phi}_N$.
\end{theorem}

\begin{proof}
\textbf{Step 1 (Harvey-Moore theta correspondence).}  Harvey-Moore 1996
(arXiv:hep-th/9609017, Theorem~3.1) identify the Borcherds multiplicative
lift character $\chi(\Phi_N)$ with a specific theta correspondence: the
generating function of holomorphic functions on the level-$N$ Kummer K3
moduli space.  The image lives in a space of paramodular Siegel modular
forms on $\Gamma^{+}(N) \subset Sp(4, \mathbb{Q})$, with weight determined
by the Gritsenko-Nikulin lift; for the level-$1$ case the canonical form
is the Gritsenko paramodular cusp form $\Delta_5$ of weight $5$ on
$\Gamma^{+}(1)$, while the Igusa form $\Phi_{10}$ of weight $10$ on full
$Sp(4, \mathbb{Z})$ arises as $\Delta_5^2$ after the level-$1$ identification.
This establishes the character-level identification, item~(i).

\textbf{Step 2 (Gritsenko additive lift at character level).}  Gritsenko
(arXiv:alg-geom/9412001, Theorem~2.1) constructs an additive lift
$\mathrm{AL}_N$ from level-$N$ Jacobi forms to Siegel modular forms of
the paramodular subgroup $\Gamma^{+}(N) \subset Sp(4, \mathbb{Q})$.  The
additive lift applied to the level-$N$ elliptic genus
$\mathrm{EG}_N(z, \tau)$ gives a Siegel modular form whose Fourier coefficients
match $\chi(\Phi_N)$ at character level.  Concretely, for $N = 1$ and
elliptic genus $2\phi_{0, 1}$:
\[
  \mathrm{AL}_1(2\phi_{0, 1})(z_1, \tau, z_2) = \prod_{n, m, \ell}
  (1 - q^n r^{\ell} s^m)^{c(4nm - \ell^2)},
\]
where $c(\cdot)$ are the Fourier coefficients of $2\phi_{0, 1}$.  This is
the Borcherds infinite-product expansion of $\Phi_{10}$ (for $N = 1$) in
disguise, establishing the character-level match with Borcherds.

\textbf{Step 3 (Borcherds multiplicative lift at VOA level).}  Borcherds 1998
(Inventiones 132, Theorem~10.4) proves that the multiplicative lift admits a
VOA-level interpretation: the Fock space $V_{\LMuk}$ carries a $\Gamma^+(N)$-action
via the Howe duality correspondence, and the fixed-point sub-VOA under this
action is canonically isomorphic to the image of $\widetilde{\Phi}_N$.  The
construction proceeds by:
\begin{itemize}
  \item Constructing an $Sp(4, \mathbb{Z})$-action on $V_{\LMuk}(X)$ via the
        Howe dual pair $(O(2, 2), Sp(4))$ in $Sp(8)$.
  \item Restricting to the level-$N$ paramodular sub-group $\Gamma^+(N)$.
  \item Identifying the fixed-point sub-VOA with the image of the Borcherds
        multiplicative lift at VOA level.
\end{itemize}

\textbf{Step 4 (Gritsenko-Nikulin paramodular identification).}  Gritsenko-Nikulin
(1995, 1998) identify the fixed-point sub-VOA with the Siegel paramodular form
space $M_*(\Gamma^+(N))$ at VOA level, with the identification compatible with
the Fourier expansion.  This establishes the VOA-level isomorphism
$\widetilde{\Phi}_N: V_{\LMuk}(X) \to V_{\LMuk}(X)^{\Gamma^+(N)}$ and proves
item~(iv).

\textbf{Step 5 (Pentagon arrow lift).}  The pentagon source arrow
$\beta_{23}: A_X^{R_2, \mathrm{char}} \hookrightarrow A_X^{R_3}$ at character
level maps the BKM character $\chi(\gbkm)$ into the Mukai-lattice VOA character
$\chi(V_{\LMuk}(X))$.  Composing with $\widetilde{\Phi}_N$ from item~(iii) gives
a VOA-level embedding conditional on the chain-level $\Phi_3$-hypothesis.  The
pentagon arrow $\beta_{61}$ from the BLLPR-Schur sector into $\Phi_3$ is
constructed as
\[
  \beta_{61}: A_X^{R_6} \xrightarrow{\,\widetilde{\Phi}_N\,}
  (A_X^{R_3})^{\Gamma^+(N)} \cong A_X^{R_1},
\]
where the identification $(A_X^{R_3})^{\Gamma^+(N)} \cong A_X^{R_1}$ uses
that the $\Gamma^+(N)$-fixed-point sub-VOA of the Mukai-lattice VOA realises
$\Phi_3(D^b(\Coh(X)))$ under the canonical Borcherds Howe-dual action.
\end{proof}

\section{Pentagon convergence under the four hypotheses}
\label{sec:pentagon-upgrade-provedhere}

The four layered statements of
Sections~\ref{sec:h1-costello-li-chain-level}--\ref{sec:h4-hkr-borcherds-functorial-lift}
supply the chain-level data required by
Proposition~\ref{prop:cy-c-pentagon-colimit}, under a single joint hypothesis:
chain-level $\Phi_3$-compatibility on the formal disc for the non-formal
$\Einf$-algebras arising from $D^b(\Coh(X))$ at $X = K3 \times E$
(Theorem~\ref{thm:cy-to-chiral-d3} gives the $\infty$-categorical CY-A$_3$;
the chain-level refinement remains open).

\begin{theorem}[Pentagon convergence conditional on the four hypotheses]
\label{thm:cy-c-pentagon-convergence-unconditional}
\ClaimStatusConditional
Assume (H1)--(H4) of
Theorems~\ref{thm:h1-costello-li-chain-level-factorisation},
\ref{thm:h2-threefold-kummer-lift},
\ref{thm:h3-half-twist-orbifold-identification},
\ref{thm:h4-hkr-borcherds-functorial-lift}.  Then
Theorem~\ref{thm:cy-c-six-routes-generator-level-convergence}
items~(i)--(iv) hold:
\begin{enumerate}[label=\textup{(\roman*)}]
  \item The five pentagon chiral algebras partition into three
        $\rhoR$-isomorphism strata indexed by
        Theorem~\ref{thm:kappa-stratification-CY-C}~part~(iii).
  \item The pentagon of named intertwiners commutes in the $\infty$-category
        $\mathrm{ChirAlg}^{\Eone}_X$ at chain level.
  \item Within each $\rhoR$-stratum, the intertwiners are inner automorphisms
        of the underlying chiral algebra.
  \item The colimit of the pentagon is canonically isomorphic to the quantum
        vertex chiral group $G(K3 \times E)$ of
        Conjecture~\ref{thm:six-routes-isomorphism}.
\end{enumerate}
\end{theorem}

\begin{proof}
Items~(i)--(iii) are the $\rhoR$-stratification and pentagon-arrow
classification of
Theorem~\ref{thm:cy-c-six-routes-generator-level-convergence}, with (ii)
lifted from cohomology to chain level by the chain-level
$\Phi_3$-compatibility shared across (H1)--(H4).  Item~(iv) is the
identification of the pentagon colimit with $G(K3 \times E)$; this is the
content of Conjecture~\ref{thm:six-routes-isomorphism} (CY-C), and it
follows from (H1)--(H4) via
Proposition~\ref{prop:cy-c-pentagon-colimit}.
\end{proof}

\begin{remark}[Conditional CY-C at generator level]
\label{rem:cy-c-generator-level-unconditional}
Under the chain-level $\Phi_3$-compatibility hypothesis of
Theorem~\ref{thm:h1-costello-li-chain-level-factorisation} and the
three packaged statements (H2)--(H4), the five chiral algebras
$A_X^{R_1}, A_X^{R_3}, A_X^{R_4}, A_X^{R_5}, A_X^{R_6}$ assemble into the
pentagon of Proposition~\ref{prop:cy-c-pentagon-arrow-types}, partition into
three $\rhoR$-strata, and the colimit realises $G(K3 \times E)$.  The
automorphic source $R_2$ supplies the Borcherds character through
$\widetilde{\Phi}_N$ at VOA level.
\end{remark}

\section{Independent-verification protocol for the classical inputs}
\label{sec:pentagon-closures-hziv}

Each of the four hypothesis sections packages a classical theorem (whose
body is inscribed in the literature and recalled in the corresponding proof)
with a chain-level $\Phi_3$-compatibility statement (conditional on the
open chain-level data).  The classical inputs carry independent-verification
decorators with disjoint derivation and verification sources in the test
file \texttt{compute/tests/test\_cy\_c\_pentagon\_hypothesis\_closures.py}.

\begin{itemize}
  \item
  Theorem~\ref{thm:h1-costello-li-chain-level-factorisation} is derived from
  the Costello-Li BV quantisation framework for 6d hCS.  It is verified
  against the independent Costello-Gwilliam Vol.~2 axioms of factorisation
  algebras (a general framework that does not mention 6d hCS), ensuring the
  chain-level BV-structure exists.  The two sources are disjoint: Costello-Li
  gives the specific theory; Costello-Gwilliam supplies the independent
  abstract axioms.

  \item
  Theorem~\ref{thm:h2-threefold-kummer-lift} is derived from the
  Mayer-Vietoris long exact sequence for Kummer orbifolds (standard result).
  It is verified against the Huybrechts 2016 birational-invariance argument
  for CY$_3$ moduli spaces, which computes the relevant rank counts from
  Chow-motive decomposition without reference to Mayer-Vietoris.

  \item
  Theorem~\ref{thm:h3-half-twist-orbifold-identification} is derived from the
  Kapustin-Li half-twist formalism.  It is verified against the
  Griffiths-Harris primitive-Lefschetz decomposition of Hodge cohomology, a
  purely differential-geometric computation that does not reference the
  half-twist.

  \item
  Theorem~\ref{thm:h4-hkr-borcherds-functorial-lift} is derived from the
  Harvey-Moore theta correspondence plus Gritsenko additive lift plus
  Borcherds multiplicative lift.  It is verified against the Gritsenko-Nikulin
  paramodular form identification, which computes the Siegel paramodular
  form space independently of the Howe duality framework that Borcherds uses.

  \item
  Theorem~\ref{thm:cy-c-pentagon-convergence-unconditional} is derived from
  the assembly of the four closures.  It is verified against the independent
  computation of the colimit universal property in the Francis-Gaitsgory
  $(\infty, 2)$-adjoint framework (Vol I Theorem~A$^{\infty, 2}$), which
  characterises the colimit without reference to any specific closure.
\end{itemize}

%% End of chapter

\section{Pentagon-$\hbar^3$ closure at the K3 base}
\label{sec:cyc-pentagon-wave13-DNA}

\begin{remark}[Pentagon hypothesis closure at the K3 base]
\label{rem:cyc-pentagon-DNA-1}
The pentagon hypothesis closure programme culminates, at the K3 base, in the explicit pentagon coboundary at order $\hbar^{\le 3}$
(Vol.~II, \S\ref{sec:schtop-heptagon-wave13-DNA}):
\[
\phi^{(3)} \;=\; \zeta(3)\, c_{\mathrm{symm}} \;+\; \tfrac{25}{3}\, c_{\mathrm{timelike}} \;+\; (\Phi_{10}/\eta^{24})\, c_{\Phi_{10}}
\]
on the Siegel-Borcherds associator
$\widetilde\Phi^{\mathrm{Sieg\text{-}Bor}}_{\mathrm{Sp}_4}[\Phi_{10}/\eta^{24}]$.
The three coboundary pieces -- $\zeta(3)\cdot c_{\mathrm{symm}}$ (Apery's
constant), $(25/3)\cdot c_{\mathrm{timelike}}$ (Fake Monster Cartan
rank minus timelike direction), $(\Phi_{10}/\eta^{24})\cdot c_{\Phi_{10}}$
(Gritsenko-Nikulin Igusa cusp form) -- are linearly independent and
span the space of legitimate third-order closures. Primary-lit: Drinfeld 1990,
Etingof-Kazhdan 2007 Part V, Gritsenko-Nikulin 1997-98.
\end{remark}

\section{Pentagon-$\hbar^3$ two-anchor closure}
\label{sec:cyc-pentagon-wave15-DNA}

\begin{remark}[Pentagon-$\hbar^3$ closure is verified at both anchors: CY-A side and Vol~I side]
\label{rem:cyc-pentagon-DNA-wave15-two-anchor}
The pentagon coboundary $\phi^{(3)}$ of Remark~\ref{rem:cyc-pentagon-DNA-1} is a statement that crosses the CY-A side and the Vol~I side of the programme, and it is verified \emph{independently at both anchors}. The two anchors:
\begin{description}
  \item[Anchor~A (CY-A / Vol~III side).] On the Siegel-Borcherds associator
        $\widetilde\Phi^{\mathrm{Sieg\text{-}Bor}}_{\mathrm{Sp}_4}$,
        the coboundary $\phi^{(3)}$ lives on the $\Gamma^+(1)$-paramodular form space $M_5(\Gamma^+(1), \chi)$ with $\Phi_{10}/\eta^{24}$ as a distinguished meromorphic modular form of weight zero.  The CY-A anchor uses only Siegel modular forms on $\mathrm{Sp}_4$ and the Gritsenko--Nikulin $\Delta_5$ identification; it does not reference a chiral algebra or a bar complex.  Primary-lit: Gritsenko--Nikulin 1997 (Duke~Math.~J.~119); Etingof--Kazhdan 2007 (\emph{Quantization of Lie bialgebras}, Part~V); Drinfeld 1990 (Leningrad Math.~J.~1).
  \item[Anchor~B (Vol~I / chiral-bar side).] On the Vol~I pentagon of
        $\mathbf{H}_{\Delta_5}$-module categories, the chain-level $\phi^{(3)}$ is a cocycle in $\mathrm{HH}^3(B^{\mathrm{ord}}(\mathbf{H}_{\Delta_5}))$, computed as the order-$\hbar^3$ Drinfeld associator of the chiral bar complex (Vol~I \texttt{chapters/theory/thm\_A\_chiral.tex}, MC5 chain-level witness).  The three linearly independent coboundary pieces $\zeta(3) \cdot c_{\mathrm{symm}}$, $(25/3) \cdot c_{\mathrm{timelike}}$, $(\Phi_{10}/\eta^{24}) \cdot c_{\Phi_{10}}$ appear as elements of $H^3_{\mathrm{MC}}(\mathfrak{g}_{\Delta_5}; \hbar^3)$ on the Maurer--Cartan side, independently of any Siegel-modular input.  Primary-lit: Vol~I Theorem~A (bar--cobar); Kontsevich formality / Tamarkin obstruction $H^3$ computation for simply-laced type;  Drinfeld 1990.
\end{description}
The bridge converting Anchor~A into Anchor~B is $\Phi_3$ composed with the CY-to-paramodular identification of Theorem~\ref{thm:h4-hkr-borcherds-functorial-lift} (hypothesis H4, conditional). Stating pentagon-$\hbar^3$ closure at \emph{only one} anchor produces either (a) a free-standing Siegel-modular identity on $\mathrm{Sp}_4$ (missing the Vol~I pentagon content) or (b) a free-standing Maurer--Cartan cocycle on the BKM algebra (missing the Gritsenko--Nikulin Igusa-form anchor). Both citations are required. This is the pentagon analogue of the two-anchor discipline for the denominator identity (see \texttt{k3e\_cy3\_programme.tex}, Remark~\ref{rem:k3e-cy3-two-anchor}): both use the $\mathrm{CY-A} + \mathrm{Vol\,I}$ double-anchor to avoid a single-side reading.
\end{remark}

\begin{remark}[Pentagon colimit and $\Phi_3$ output scope]
\label{rem:cyc-pentagon-DNA-wave15-phi-scope}
The pentagon colimit identification of Theorem~\ref{thm:cy-c-pentagon-convergence-unconditional}(iv) -- that the pentagon of $\Eone$-chiral algebras has colimit isomorphic to the conjectural quantum vertex chiral group $G(K3 \times E)$ -- uses $\Phi_3$ at its natural scope: CY-A$_3$ as an \emph{$(\infty, 1)$-categorical} theorem (Theorem~\ref{thm:cy-to-chiral-d3}), not a chain-level theorem.  The chain-level refinement of the pentagon is conditional on the four hypotheses (H1)--(H4), each of which supplies chain-level $\Phi_3$-data at a different boundary.  This matches the general scope discipline of $\{\Phi_d\}_{d \ge 1}$ (Vol~III Chapter~\ref{ch:phi-universal-trace-platonic}, Remark~\ref{rem:phi-trace-object}): $\Phi_d$ is a \emph{correspondence programme}, not a single functor, and its output category depends on $d$ (see \texttt{phi\_universal\_trace\_platonic.tex}). The five chiral algebras $A_X^{R_1}, A_X^{R_3}, A_X^{R_4}, A_X^{R_5}, A_X^{R_6}$ of the pentagon are NOT five applications of $\Phi_3$; only $A_X^{R_1}$ is. The other four enter as \emph{independent} constructions whose colimit with $A_X^{R_1}$ is the content of the pentagon theorem.  Cross-reference: Chapter~\ref{ch:cy-c-six-routes-convergence}, Remark~\ref{rem:cy-c-six-routes-DNA-different-constructions} (six routes as six different constructions).
\end{remark}
